% ============================================================
%  8. CONCLUSIONES Y RECOMENDACIONES   (guía, 2.8)
%  Regla explícita: las conclusiones deben derivarse de los
%  resultados y responder al objetivo general y a los específicos.
%  NO deben introducir información nueva no desarrollada antes.
% ============================================================
\chapter{Conclusiones y Recomendaciones}
\label{cap:conclusiones}

\section{Conclusiones}
\label{sec:conclusiones}

\subsection{Conclusión respecto del objetivo general}

La pregunta de investigación --en qué zonas del eje metropolitano de Cochabamba se
concentra la demanda de información de transporte no resuelta y qué características del
territorio se asocian con esa concentración (sección~\ref{sec:formulacion-problema})-- se
responde con evidencia convergente de tres análisis independientes
(sección~\ref{sec:interpretacion-resultados}): la brecha de cobertura centro-periferia es
real y estadísticamente significativa (H1, $p=6{,}1\times10^{-9}$); el gradiente espacial de
demanda es real y el modelo lo reproduce ($\rho=-0{,}615$ entre distancia al centro y
demanda observada); y el valor predictivo del modelo está concentrado exactamente en las
celdas de mayor demanda (H2), no distribuido de manera uniforme. El objetivo general
--construir un diagnóstico espacial de las brechas de información de transporte público y
derivar de él una lista priorizada de zonas de mapeo (sección~\ref{sec:objetivo-general})--
se cumple: el proyecto entrega tanto el diagnóstico (indicadores territoriales,
sección~\ref{sec:tabla-indicadores}, y contraste formal de la brecha, H1) como la lista
priorizada accionable (top-20 de celdas, sección~\ref{sec:priorizacion}), con evidencia de
que esta última no depende del modelo predictivo para ser útil.

A continuación se presentan las conclusiones correspondientes a cada objetivo específico.

\begin{enumerate}
  \item [OE1.] \textbf{Consolidar y auditar los archivos de consultas.} Se cumple
        íntegramente. Se consolidaron 1.927.675 registros de 85 exportaciones semanales sin
        pérdida ni duplicación silenciosa de filas (sección~\ref{sec:auditoria-esquema}),
        con las anomalías de esquema y codificación de seis archivos resueltas y
        documentadas, y la naturaleza de la variable \texttt{distancia} verificada
        empíricamente en lugar de asumida (sección~\ref{sec:validacion-distancia}). El
        resultado es un conjunto base donde cada fila excluida tiene un motivo registrado
        (tabla~\ref{tab:flujo-filtros}), cumpliendo el objetivo de trazabilidad completa
        entre el dato crudo y el dato modelado.
  \item [OE2.] \textbf{Caracterizar la demanda mediante análisis exploratorio
        espacio-temporal.} Se cumple. La demanda está fuertemente concentrada en el espacio
        (coeficiente de Gini de \textbf{0,934} en resolución 8, sección~\ref{sec:sensibilidad-h3})
        y decae con la distancia al centro; ambos patrones son estables ante cambios de
        resolución H3 --correlaciones de rangos superiores a 0,98 entre resoluciones
        adyacentes-- lo que descarta que sean un artefacto de la escala de agregación
        elegida (problema de la unidad de área modificable, capítulo~\ref{cap:marco-teorico}).
        En el tiempo, la demanda no es estacionaria y presenta un vacío estructural de siete
        semanas que condicionó la partición de entrenamiento y prueba
        (sección~\ref{sec:particion-cronologica}).
  \item [OE3.] \textbf{Construir indicadores territoriales por celda y semana.} Se cumple.
        La tabla de indicadores (sección~\ref{sec:tabla-indicadores}) integra demanda,
        cobertura GTFS, centralidad y composición temporal sobre 42.676 observaciones
        celda-semana, con el método de cálculo de cada variable documentado y nueve
        variables excluidas explícitamente por fuga de información
        (sección~\ref{sec:seleccion-algoritmos}). No se incorporaron Puntos de Interés de
        OpenStreetMap como variable de entorno urbano, decisión justificada por su
        cobertura desigual (capítulo~\ref{cap:alcance}) y retomada como recomendación más
        abajo.
  \item [OE4.] \textbf{Contrastar formalmente las hipótesis de brecha y de mejora
        predictiva.} Se cumple, con dos resultados que se refuerzan mutuamente. H1 se
        confirma: las celdas periféricas tienen una tasa de demanda no resuelta
        significativamente mayor que las centrales (Mann-Whitney,
        $p=6{,}1\times10^{-9}$, sección~\ref{sec:contraste-hipotesis}). H2 se confirma con
        matices: el test de Diebold-Mariano sobre la serie semanal no alcanza significancia
        por falta de potencia (8 puntos), pero el Wilcoxon pareado por celda sí
        ($p=0{,}033$), y la ventaja de Random Forest está concentrada en las celdas de
        mayor demanda --el segmento operativamente más relevante-- y no distribuida de
        forma uniforme.
  \item [OE5.] \textbf{Predecir demanda con modelos comparados contra líneas base.} Se
        cumple. Se compararon cuatro modelos de tres familias (Ridge, Lasso, Random Forest,
        XGBoost) bajo validación por ventanas deslizantes, nunca aleatoria
        (sección~\ref{sec:proceso-entrenamiento}). El hallazgo más informativo no fue qué
        modelo ``ganó'' en abstracto sino por qué los modelos lineales fallan: Ridge y
        Lasso extrapolan de forma inestable ante la tendencia de crecimiento sostenido,
        produciendo R\textsuperscript{2} fuertemente negativo en prueba pese a un ajuste
        razonable en entrenamiento --evidencia de que la relación entre demanda y
        condiciones territoriales no es lineal. Random Forest se seleccionó como modelo
        final con criterios explícitos, no por haber sido el primero ejecutado
        (sección~\ref{sec:seleccion-final}): MAE de prueba de \textbf{10,19} (o
        \textbf{5,91} excluyendo dos semanas de prueba que resultaron ser fragmentos de un
        solo día, sección~\ref{sec:semanas-parciales}), la evaluación más conservadora y la
        más representativa reportadas juntas, como corresponde cuando ninguna de las dos
        cifras por sí sola sería honesta.
  \item [OE6.] \textbf{Derivar una regla de priorización y contrastar el aporte del
        modelo.} Se cumple, con un resultado negativo declarado como tal. La regla de
        decisión --demanda esperada multiplicada por tasa de no cobertura-- produce un
        top-20 de celdas que concentra el 70,5~\% de la demanda no resuelta estimada
        (sección~\ref{sec:priorizacion}), y esas celdas no son las de mayor demanda
        absoluta sino zonas de demanda baja o media mal cubiertas --el mismo hallazgo de H1
        visto desde la decisión operativa. Pero el criterio preinscrito que evaluaría si el
        modelo predictivo mejora ese ordenamiento frente a una media móvil de 4 semanas
        \textbf{no se cumplió} (sección~\ref{sec:sirve-para-priorizar}): Random Forest no
        supera a la media móvil en ninguna de las seis semanas completas de prueba, en
        ninguna métrica de ordenamiento. Este resultado negativo, obtenido con un criterio
        declarado antes de observar los datos, es evidencia de rigor y no una falencia: para
        decidir qué celdas mapear a un horizonte de una semana no se justifica un modelo
        complejo, y la contribución del proyecto está en el pipeline reproducible, la
        caracterización territorial y el diagnóstico de cobertura, no en la superioridad de
        un estimador.
  \item [OE7.] \textbf{Proponer una arquitectura de actualización periódica.} Se cumple. Se
        construyó y se ejerció --no solo se documentó-- un prototipo real
        (\var{src/trufi_ds/api.py}) que sirve predicciones reconstruyendo las mismas
        variables usadas en entrenamiento, con ejemplos de consumo obtenidos de ejecuciones
        reales (sección~\ref{sec:despliegue}). Ejercitarlo expuso una limitación operativa
        real, no hipotética: la última semana del dataset es en sí misma una semana
        parcial, por lo que la primera predicción en producción heredaría la misma
        contaminación de rezago detectada en la evaluación. El plan de monitoreo deriva sus
        umbrales de la distribución real de error y volumen de este mismo proyecto,
        incluyendo un chequeo de completitud de datos que existe porque este proyecto
        encontró el problema que previene.
\end{enumerate}

\subsection{Limitaciones reconocidas}

Las limitaciones declaradas de antemano en el capítulo~\ref{cap:alcance} se sostienen, y se
completan con las que emergieron durante el desarrollo:

\begin{itemize}
  \item El feed GTFS usado para calcular cobertura tiene un período de validez declarado
        (2024-2026) posterior a buena parte de las consultas analizadas (desde 2022): la
        cobertura se mide contra el estado vigente del mapeo, no contra el estado que tenía
        en cada momento histórico.
  \item La cobertura se mide contra el trazado de la ruta (\texttt{shapes.txt}) y no contra
        las paradas (\texttt{stops.txt}), lo que mide accesibilidad a la línea de transporte
        y no al punto de abordaje efectivo, y por tanto sobreestima la cobertura tal como la
        percibe el usuario.
  \item La partición central/periferia usada en H1 se basa en la mediana de distancia al
        centro (15,49~km), un corte simple y documentado pero arbitrario, que no
        corresponde a ningún límite administrativo real.
  \item El test de Diebold-Mariano de H2 tiene solo ocho puntos, potencia insuficiente por
        construcción; su resultado no significativo no implica equivalencia entre modelos,
        y se complementó con el Wilcoxon pareado por celda por esa misma razón.
  \item Una celda ``sin cobertura'' puede tener servicio de transporte no mapeado en el
        GTFS de Trufi: la métrica localiza la brecha de \emph{información}, que es el
        objeto de este estudio, y no debe leerse como ausencia real de transporte.
  \item La demanda hereda el sesgo de adopción de la aplicación: zonas con menos usuarios de
        Trufi generan menos consultas aunque tengan mayor necesidad de transporte.
  \item El identificador de usuario es de nivel instalación, no de persona; una
        reinstalación genera un identificador nuevo y un dispositivo compartido produce el
        caso inverso.
  \item No se incorporaron Puntos de Interés de OpenStreetMap como variable de entorno
        urbano, ni un componente de segmentación por agrupamiento, por las razones
        declaradas en el capítulo~\ref{cap:alcance}.
  \item No existe una suite de pruebas automatizadas sobre el pipeline ni sobre las métricas
        del modelo, reconocida como la principal deuda técnica del proyecto.
\end{itemize}

Estas limitaciones acotan el alcance de las conclusiones anteriores sin invalidarlas: cada
hallazgo reportado --la brecha de cobertura, el gradiente espacial, la ventaja concentrada
del modelo, la refutación de su aporte al ordenamiento-- se sostiene dentro de estos límites
declarados, no a pesar de ocultarlos.

\section{Recomendaciones}
\label{sec:recomendaciones}

Las recomendaciones se derivan de los resultados obtenidos y de las limitaciones
identificadas; ninguna es genérica.

\begin{enumerate}
  \item [R1.] \textbf{Completar (\textit{backfill}) las semanas de datos parciales antes de
        cualquier despliegue real}: 2024-W18 y 2024-W23 con datos reales de esas semanas
        calendario completas (sección~\ref{sec:despliegue}); de lo contrario, la primera
        predicción en producción hereda el mismo artefacto que ya distorsionó la evaluación.
  \item [R2.] \textbf{Formalizar el chequeo de completitud de datos} (piso de 10.041
        consultas semanales, sección~\ref{sec:monitoreo}) como validación automática dentro
        del pipeline de actualización, no solo como umbral documentado --habría detectado
        el artefacto de semanas parciales en el momento en que ocurrió.
  \item [R3.] \textbf{Revisar la inestabilidad de extrapolación de los modelos lineales}
        antes de usarlos como referencia interpretable en producción: opciones concretas
        incluyen regresión cuantílica o un techo de extrapolación aprendido en vez del
        umbral fijo usado en este trabajo.
  \item [R4.] \textbf{Ampliar la validación de H1} más allá de la partición por mediana de
        distancia al centro, usando límites administrativos reales de periferia si se
        consigue una fuente confiable, en vez de un corte simple pero arbitrario.
  \item [R5.] \textbf{Aumentar la ventana de prueba} en un futuro ciclo de recolección de
        datos para dar más potencia al test de Diebold-Mariano, sin sacrificar datos de
        entrenamiento.
  \item [R6.] \textbf{Evaluar el aporte del modelo a horizontes de predicción mayores a una
        semana}: la refutación de la sección~\ref{sec:sirve-para-priorizar} vale
        específicamente para $h=1$ semana; a horizontes mayores la media móvil envejece y
        el modelo podría mostrar ventaja, pero esa prueba requiere reconstruir los rezagos
        a $h=2,3,4$ semanas y reentrenar, y queda pendiente.
  \item [R7.] \textbf{Incorporar Puntos de Interés de OpenStreetMap y un componente de
        segmentación territorial por agrupamiento} como extensiones del perfil de celda,
        considerados en el diseño original de este proyecto y descartados por razones de
        cobertura de datos y de alcance (capítulo~\ref{cap:alcance}), no por falta de
        valor potencial.
  \item [R8.] \textbf{Confirmar con el equipo de backend de Trufi App} si el cambio de
        esquema y codificación detectado en seis archivos, justo después del vacío de siete
        semanas, corresponde a un cambio real del proceso de exportación
        (sección~\ref{sec:auditoria-esquema}) --queda como pregunta abierta, nunca
        confirmada con la fuente.
  \item [R9.] \textbf{Antes de un despliegue productivo real}: agregar autenticación,
        límites de tasa, un orquestador para el reentrenamiento trimestral en vez de
        ejecución manual de scripts, y pruebas automatizadas de regresión sobre las
        métricas del modelo (sección~\ref{sec:despliegue}).
\end{enumerate}
